
\section{Additional mathematical notes}
\label{app:additional_notes}

\subsection{On unique ergodicity of Kronecker rotations}

For $\alpha$ satisfying \eqref{eq:rational_independence}, the translation $x\mapsto x+\alpha$ on $\TT^d$ is uniquely ergodic with respect to Haar measure. As a consequence, for continuous $f$, Birkhoff averages converge uniformly in the initial condition. In the main text we used the equivalent uniform distribution formulation because it interfaces directly with discrepancy bounds and Riemann integrable kernels.

\subsection{On discrepancy estimates and Diophantine properties}

The star discrepancy of the Kronecker sequence $\{x_0+t\alpha\}$ admits explicit bounds in terms of Diophantine approximation properties of $\alpha$; in $d=1$ these are controlled by continued fractions \cite{KuipersNiederreiter1974,IosifescuKraaikamp2002}. This paper does not require a specific bound, only that discrepancy provides an auditable sampling error term (Corollary~\ref{cor:sampling_error}). The selection principle (Section~\ref{sec:selection_principle}) uses these estimates as motivation.

\subsection{A closed 1D rate bound for constant-type rotations (bounded partial quotients)}
\label{app:1d_discrepancy_constant_type}

For one-dimensional scans $x_t=x_0+t\alpha\pmod{1}$, finite-horizon stability can be made fully explicit under a standard Diophantine regularity condition.

\begin{proposition}[Star discrepancy for bounded partial quotients]
\label{prop:1d_discrepancy_log_over_n}
Let $\alpha=[0;a_1,a_2,\dots]$ be irrational with bounded partial quotients $a_n\le A$. Then the star discrepancy of the Kronecker point set $P_N=\{\{x_0+t\alpha\}:0\le t\le N-1\}\subset[0,1)$ satisfies
\begin{equation}
\label{eq:1d_discrepancy_log_over_n}
D_N^\ast(P_N) \le \frac{2A(3+\log_\varphi N)}{N},
\end{equation}
where $\varphi=(1+\sqrt{5})/2$.
\end{proposition}

\begin{proof}
Let $q_n$ denote the continued-fraction convergent denominators of $\alpha$, and write the Ostrowski expansion $N=\sum_{j=0}^m b_j q_j$. For $u\in[0,1]$, set $f_u(x)=\ind_{[0,u)}(x)-u$, so that $\int_0^1 f_u=0$ and $\Var(f_u)\le 2$.

By the Denjoy--Koksma inequality at convergent lengths \cite{Denjoy1932,Koksma1935,KatokHasselblatt1995}, each block sum over a length-$q_j$ orbit segment is bounded by $\Var(f_u)$ uniformly in the starting point. Applying an Ostrowski block decomposition and summing over the $b_j$ blocks yields
\[
\left|\sum_{t=0}^{N-1} f_u(x_0+t\alpha)\right|\le \Var(f_u)\sum_{j=0}^m b_j \le 2\sum_{j=0}^m b_j.
\]
Dividing by $N$ and taking the supremum over $u$ gives
\[
D_N^\ast(P_N)\le \frac{2}{N}\sum_{j=0}^m b_j.
\]
If $a_n\le A$, then Ostrowski digits satisfy $b_j\le a_{j+1}\le A$, hence $\sum_{j=0}^m b_j\le A(m+1)$.
Moreover, $q_{n+1}=a_{n+1}q_n+q_{n-1}\ge q_n+q_{n-1}$ implies $q_n\ge F_n$ for Fibonacci numbers, and $F_n\ge \varphi^{n-2}$ for $n\ge 2$. Since $q_m\le N$, we obtain $\varphi^{m-2}\le N$, hence $m\le 2+\log_\varphi N$. Therefore $m+1\le 3+\log_\varphi N$ and
\[
D_N^\ast(P_N)\le \frac{2A(3+\log_\varphi N)}{N}.
\]
\end{proof}

\begin{remark}
For the golden branch $\alpha=\varphi^{-1}$ one has $A=1$, giving an explicit certified $O((\log N)/N)$ discrepancy bound. Combined with Proposition~\ref{prop:explicit_1d_variations}, this turns the $\log 2$ and $\pi$ scan realizations into fully quantified, parameter-free finite-$N$ bounds.
\end{remark}


